% Generated draft: requires mathematical and source review.
\documentclass[11pt]{article}
\usepackage[margin=25mm]{geometry}
\usepackage{fontspec}
\setmainfont{Latin Modern Roman}
\usepackage{amsmath,amssymb,mathtools,unicode-math}
\setmathfont{Latin Modern Math}
\usepackage{enumitem,needspace,xurl,hyperref}
\setlength{\parindent}{0pt}
\setlength{\parskip}{5pt}
\setlength{\emergencystretch}{3em}
\newcommand{\N}{\mathbb{N}}
\newcommand{\Z}{\mathbb{Z}}
\newcommand{\Q}{\mathbb{Q}}
\newcommand{\R}{\mathbb{R}}
\newcommand{\C}{\mathbb{C}}
\newcommand{\F}{\mathbb{F}}
\newcommand{\A}{\mathbb{A}}
\newcommand{\PP}{\mathbb{P}}
\newcommand{\E}{\mathbb{E}}
\newcommand{\eps}{\varepsilon}
\newcommand{\dd}{\,\mathrm{d}}
\begin{document}
\section*{UnsolvedMath ID: 30000973}
\textbf{Source problem code:} catalogue entry
\subsection{Definitions and mathematical statement}
For an intended number of markings \(n\), let \(\mathcal M_{1,n}\) denote the moduli of smooth connected genus-one curves with \(n\) distinct labelled marked points. The source names alternate birational models \(M_{1,n}(m)\) and \(\widetilde M_{1,n}(m)\), with the latter described as obtained from normalization of a reduced stack. It does not give the stability condition, the stack, its coarse-space construction, the parameter range, or a ground field.

Fix \(n\), and write \(I_n\) for the source's unstated admissible values of \(m\). The requested equality is, geometrically, the canonical identification
\[
 \forall m\in I_n,\qquad
 M_{1,n}(m)\ \cong\ \widetilde M_{1,n}(m),
\]
extending the identification on the common smooth pointed-elliptic-curve locus. The record says every \(m\), rather than only small values or an unspecified sufficiently large range. It does not ask merely for birational equivalence, which is already part of their description as models of the same moduli problem. Literal set equality would depend on a chosen construction and is not the invariant mathematical target.

For an object \(\mathfrak X\), reduction removes nilpotents in its structure sheaf, and normalization, when defined in the relevant category, gives a morphism
\[
 \nu:(\mathfrak X_{\mathrm{red}})^\nu\longrightarrow\mathfrak X_{\mathrm{red}}.
\]
If the two named objects were precisely a reduced object and its normalization, the requested identification would amount to this normalization morphism being an isomorphism. The extraction does not identify them sufficiently to make that equivalence unconditional, especially because stacks and coarse moduli spaces may be involved.

The supplied assessment dated 2026-08-21 calls the equality one between an \(m\)-stable model and a \(\operatorname{Proj}\) model, while the original and cleaned statements describe normalization of a reduced stack. No graded ring defining the \(\operatorname{Proj}\) model, or theorem relating it to that normalization, is supplied. These descriptions must be reconciled; they are not silently assumed identical. The assessment reports the all-\(m\) equality as unresolved, but supplies no complete formulation or independent later-source details.
\subsection{Short English statement}
The question asks whether two birational compactification models of pointed genus-one curves agree for every allowed stability parameter. The constructions and parameter range are absent, and the supplied normalization description must be reconciled with the assessment's description as a Proj model.
\subsection{Sources}
\begin{itemize}
\item Local UnsolvedMath registry, permanent ID \(30000973\), exact source code \texttt{catalogue entry}, proposed year 2008; no individual proposer is supplied.
\item \emph{Classical Algebraic Geometry (2009)}, Oberwolfach Reports, DOI 10.4171/owr/2008/26, pp. 1425--1474. \url{https://doi.org/10.4171/owr/2008/26}.
\item Supplied source labelled ``report text recording the unresolved equality'': \url{https://publications.mfo.de/bitstream/handle/mfo/3070/OWR_2008_26.pdf?isAllowed=y&sequence=1}.
\item The supplied \texttt{corrected\_verified} note reports editorial and LaTeX cleanup without added source-dependent content. The open label and assessment checked 2026-08-21 report an unresolved all-\(m\) equality between \(m\)-stable and \(\operatorname{Proj}\) models. This is dated supplied triage; its descriptions and claim about later sources require independent verification.
\end{itemize}
\end{document}
